\documentclass{llncs}

\usepackage[utf8]{inputenc}

\begin{document}

\title{Metodologías Ágiles, Scrum y Certificaciones}
\author{Nahuel Toci}
\institute{Universidad Nacional de Cuyo}

\maketitle

\begin{abstract}
Este informe resume los conceptos fundamentales de las metodologías ágiles, el marco de trabajo Scrum y los niveles de certificación White Belt y Green Belt, presentados de forma sintética para su estudio.
\end{abstract}

\section{Introducción}
Las metodologías ágiles surgen como alternativa a los modelos tradicionales, permitiendo adaptarse a cambios y mejorar la entrega de valor al cliente mediante iteraciones cortas.

\section{Metodologías Ágiles}
Son enfoques de desarrollo basados en iteración, colaboración y flexibilidad.

\subsection{Características principales}
\begin{itemize}
\item Desarrollo iterativo e incremental
\item Entregas frecuentes de valor
\item Adaptación al cambio
\item Trabajo en equipo
\end{itemize}

\subsection{Diferencia con modelo tradicional}
\begin{itemize}
\item Tradicional: planificación rígida (Waterfall)
\item Ágil: adaptación continua
\end{itemize}

\section{Scrum}
Scrum es un marco ágil para gestionar proyectos complejos mediante ciclos llamados Sprints.

\subsection{Roles}
\begin{itemize}
\item Product Owner: define requerimientos (Product Backlog)
\item Scrum Master: facilita el proceso y elimina impedimentos
\item Equipo de desarrollo: construye el producto
\end{itemize}

\subsection{Eventos}
\begin{itemize}
\item Sprint: ciclo de trabajo (1-4 semanas)
\item Daily Scrum: reunión diaria breve
\item Sprint Review: revisión del incremento
\item Sprint Retrospective: mejora continua del equipo
\end{itemize}

\subsection{Artefactos}
\begin{itemize}
\item Product Backlog: lista de requisitos
\item Sprint Backlog: tareas del Sprint
\item Incremento: resultado funcional
\end{itemize}

\section{White Belt y Green Belt}
Certificaciones asociadas a mejora de procesos (Lean / Six Sigma).

\subsection{White Belt}
\begin{itemize}
\item Nivel básico
\item Conocimiento general
\item Participación en proyectos
\end{itemize}

\subsection{Green Belt}
\begin{itemize}
\item Nivel intermedio
\item Lidera proyectos pequeños
\item Aplica metodologías como DMAIC
\end{itemize}

\section{Relación con metodologías ágiles}
Las certificaciones complementan a Scrum aportando herramientas de mejora continua y optimización de procesos.

\section{Conclusión}
Las metodologías ágiles, especialmente Scrum, permiten mejorar la eficiencia, adaptabilidad y calidad en proyectos. Las certificaciones como Green Belt refuerzan el enfoque de mejora continua.

\end{document}